% Modernised reading — fonds Alexandre Grothendieck, shelfmark 128
% (pages 1-31, the whole folder).
%
% This is an INTERPRETATION, not a transcription. It re-reads the pages in
% today's notation and vocabulary, states the hypotheses the manuscript leaves
% implicit, and drops the critical apparatus so that the mathematics reads
% straight through. Everything the brackets and underlines carried moves into a
% footnote. It is held to being mathematically correct as it stands; where the
% manuscript is loose or mistaken, the true statement is what appears here and
% the footnote says what the page has. In any dispute of substance the
% transcription governs: whoever cites Grothendieck must cite batch-01.fr.tex
% (pp. 1-20) or batch-02.fr.tex (pp. 21-31).
%
% Pass: Opus 5.5 (claude-opus-5-5), 2026-10-03 — first pass, unchecked against the pages by a human.
% The folder was read whole, its two batches together; this is one document
% for the shelfmark. It was made on Opus 5.5 and has not been measured against
% the reference reading of folder 115 (made on Opus 5). The literature named
% in the footnotes is cited from memory and was not checked against the
% sources.
%
% Scope. Pages 1-29 are his. Page 30 is a blank cover sheet. Page 31 is a
% sheet in another hand, in English; it is not restated beyond what the
% transcription's note says (RIGHTS.md). The « Notes Murre-Levelt » of the
% inventory title are not otherwise present.
%
% Spine. Three runs, one method: compare the Amitsur complex of A'/A with that
% of its reduction modulo a nilpotent (or square-zero) ideal.
%   pp. 1-11   « Réductions via le conducteur » (his circled 1, « N° 1 »):
%              I contained in the conductor of A in A'; exact sequence of
%              cosimplicial modules; effective descent for flat modules when
%              I is nilpotent (Prop. 2 and Cor. 1-3); augmentation ideals and
%              equivalence relations transported mod I (Prop. 3).
%   pp. 12-18  « Résultats auxiliaires sur les relations d'équivalence »
%              (circled 2): extensions with A'/A = k (Prop. 1); a criterion
%              for length(A'/A) <= 1 (Prop. 2), stopped at « il faut prouver
%              que V est de dim. 1 ».
%   pp. 19-29  infinitesimal neighbourhoods: Lemme 1 (G_m vs G_a mod a
%              square-zero ideal) and its corollaries for Gl_n, Aut(E);
%              Lemme 2 (making a matrix invertible); the Théorème on A
%              semi-local noetherian, A_n = A/m^{n+1}, and its proof by
%              induction on n and Mittag-Leffler.
%
% Notation fixed for the document.
%  - A' is the A-algebra throughout. Pages 1-11 call it B (and write B = A'
%    on p. 4); pages 19-29 use B for H^0(A'/A). Here B = H^0(A'/A) only.
%  - C^n(A'/A) = A'^{⊗_A(n+1)}, the Amitsur complex C^•(A'/A); the page writes
%    C^n, K(B/A) (pp. 1-3) and a doubled C with a dot (pp. 19-29).
%    H^i(A'/A, M) its cohomology with coefficients. « Simplicial » on the
%    page is « cosimplicial » here (footnoted once, p. 2).
%  - \widetilde{B}_0 = H^0(A'_0/A_0), the page's C_0 (pp. 19-22); the ring of
%    invariants of an equivalence ideal J (p. 11's C) is A'^J.
%  - Gothic m is the radical (pp. 25-29) or the maximal ideal (pp. 12-18); the
%    integer of the Théorème's (ii_m) is written r here (the page writes m).
%
% Substantive departures, each footnoted where it occurs:
%  p. 4    « H^0(B, A) » read H^0(A'/A); footnote (*): being a descent
%          morphism for flat modules is equivalent to A = H^0(A'/A) alone;
%          the hypothesis H^1(A'/A) = 0 of Prop. 2 is not used by the proof.
%  p. 5    « A'_0/A' » read A'_0/A_0.
%  p. 6    the illegible end of the I^2 = 0 case reconstructed (ours):
%          Nakayama plus flatness of E'.
%  p. 6    flatness of E by the local flatness criterion for a nilpotent ideal
%          (named, ours).
%  p. 9    « Y un A-préschéma quelconque » is false: Cor. 3 needs Y flat; a
%          counterexample (A = k[u,v]/(u,v)^2 in (k[t]/t^2)^2) is given (ours).
%  p. 13   Prop. 1 (iii) uses N° 1 Prop. 2, which needs m nilpotent; the
%          reading states (iii) for m nilpotent and flat base change; the
%          page's general claim is left open. Case k[ε] of (i) verified (ours).
%  p. 14   the margin's « extensions résiduelles triviales » is made part of
%          Prop. 2; « J nilpotent » justified from it (ours).
%  p. 15   « δ_0(x'_0) ∈ A'_0 » read ∈ Δ_0 ⊂ A''_0.
%  p. 18   « multiples de δ(w) » read δ(v); dim V = 1 proved (ours).
%  p. 20   the square of p. 20 does not commute as drawn: the multiplicative
%          connecting map is u ↦ u^{-1} ∂(u), not ∂(u). Lemme 1 is stated as
%          his claim, with the gap said in the body; same for Cor. 2-3
%          (pp. 21-22) and for the steps of the Théorème that use them.
%  p. 21   the doubtful « Z C^•(A'/A, M_n) » read I C^•(A'/A, M_n), following
%          p. 19 and batch-01's header.
%  p. 23   Lemme 2's hypothesis (i) is implied by (ii) (finite fields are
%          perfect) (ours).
%  p. 25   « ½ local noethérien » over a struck « artinien » kept.
%  p. 27   descent data relative to A'_n/A_n and to A'_n/B^(n) coincide (his
%          « i.e. », justified, ours).
%  p. 28   Mittag-Leffler for H^0 supplied (finite length, ours).
%  p. 29   the « lemme facile de passage à la limite » is not written; what it
%          needs (Artin-Rees) said, not proved.
%
% Announced and never established: the effectivity criterion of p. 11 (c)
% (its « … ] » is unfinished; checked here, ours); case (ii) of p. 13; the end
% of p. 18 (supplied, ours); the commutativity behind Lemme 1 (p. 20); most of
% the proof of Lemme 2 (p. 23) and of its variant (p. 24, nearly illegible);
% the N.B. of p. 24 on a residue extension k(t); the diagonal N.B. of p. 25;
% large parts of the proof of a) (pp. 26-27, illegible); the deduction of
% b)(i) from (iii) « en faisant m = 1, m = 2 » (p. 28); the limit lemma
% (p. 29); the whole of c) (p. 29 stops in a bracket).
\documentclass[11pt,a4paper]{article}
\input{../preamble/grothendieck}

\folder{128}
\pages{1}{31}
\dating{s.d. — le groupe « Descentes » (126 à 133) est daté [avant 1970]}
\watermark{Édition de démonstration\\interprétation personnelle de l'œuvre}
\foldertitle{Descente non plate. Modules formels. Cf Notes Murre-Levelt : notes manuscrites (s.d.) — lecture modernisée du dossier entier}

\begin{document}

\section*{Résumé}

\begin{resume}
Ces feuillets cherchent à quelles conditions on peut reconstituer un objet
algébrique à partir de ce qu'on en voit sur un « revêtement », quand ce
revêtement est mal fait.

Le problème est celui de la descente. Pour construire un module sur un anneau
$A$, il est souvent plus facile de le construire sur un anneau plus gros $A'$,
puis de dire comment les morceaux se recollent : c'est ce qu'on appelle une
donnée de descente. On recolle ainsi des cartes pour faire un atlas : chaque
carte est une donnée locale, et la règle de recollement dit comment deux
cartes se raccordent sur leur partie commune. La théorie classique, celle
de Grothendieck lui-même, marche très bien quand $A'$ est \emph{plat} sur $A$ — en gros, quand $A'$ ne crée ni ne
détruit rien en se superposant à $A$. Le titre du dossier dit d'emblée de quoi
il s'agit : la descente \emph{non plate}.

L'exemple qu'il faut garder en tête est celui d'une courbe qui se recoupe en
un point, comme un 8 aplati, qu'on « déplie » en une droite : la droite
revêt la courbe, mais deux de ses points tombent sur le même point de la
courbe. Les fonctions sur la courbe sont exactement les fonctions $f$ sur la
droite telles que $f(0) = f(1)$. Les deux anneaux, celui de la droite et
celui de la courbe, coïncident sur un gros morceau — les fonctions nulles en
$0$ et en $1$ — et ne diffèrent que par ce qui se passe au point double. Ce
morceau commun est le \emph{conducteur}. La première idée du dossier est
d'exploiter cette coïncidence : tout ce qui se passe dans le conducteur est
le même en haut et en bas, de sorte que le problème de descente se ramène à
celui qu'on obtient en coupant le conducteur, souvent beaucoup plus simple —
plat, par exemple. Quand le conducteur est fait d'éléments nilpotents, cette
réduction suffit à prouver que la descente marche.

La deuxième idée est celle des approximations successives, qu'on connaît par
les développements limités : on résout un problème à l'ordre $0$, puis on
corrige à l'ordre $1$, puis à l'ordre $2$, et ainsi de suite, chaque
correction étant un problème linéaire. Grothendieck l'applique à un anneau
local complet, qu'il remplace par ses quotients successifs
$A/\mathfrak{m}^{n+1}$ : à chaque étape, l'obstruction à recoller des
modules libres de rang $1$ (un problème multiplicatif) se compare à
l'obstruction pour un problème additif, beaucoup plus simple, puis on passe à
la limite. C'est le sens du second mot du titre, « modules formels ».

Entre les deux, sept pages étudient les \emph{relations d'équivalence} qu'on
peut mettre sur un tel revêtement — les façons d'y identifier des points — et
montrent qu'une extension où ne s'ajoute « qu'un seul élément de plus »
n'en admet que deux, la plus fine et la plus grossière.

Ce sont des notes de travail au sens strict. Les énoncés se corrigent en
cours de route — un « artinien » rayé devient « semi-local noethérien », une
hypothèse ajoutée dans la marge en fait silencieusement partie —, la dernière
démonstration s'arrête au milieu d'une parenthèse, et une page s'interrompt
sur « il faut prouver que ». Deux points de ses arguments ne passent pas tels
quels, et on le dit à leur place ; on y voit aussi comment il avance : réduire
d'abord au cas où la base est un corps, puis regarder ce qu'il reste. La
dernière feuille, d'une autre main et en anglais, appartient sans doute aux
notes de Murre et Levelt que cite le titre ; rien sur la page ne le dit.

Les noms sous lesquels chercher la suite : complexe d'Amitsur, morphisme de
descente effective, carré de Milnor et pincement, critère local de
platitude, relations d'équivalence finies et quotients, cohomologie non
abélienne, condition de Mittag-Leffler.

\keywords{descent theory, effective descent morphism, non-flat descent, Amitsur complex, cosimplicial module, conductor ideal, Milnor square, pinching, local flatness criterion, finite equivalence relation, algebra of dual numbers, nonabelian cohomology, logarithmic derivative, Mittag-Leffler condition, formal completion}
\end{resume}

\pagerange{1}{31}
\section*{Le fil du dossier, et les conventions}

\subsection*{Les stations}

L'inventaire range le dossier dans le groupe « Descentes » (cotes 126 à 133),
daté « [avant 1970] » ; lui-même est « s.d. », et rien sur les feuillets ne le
date. Trois suites, que Grothendieck numérote en partie lui-même :

\begin{itemize}
  \item[I.] \emph{Pages 1 à 11}, « Réductions via le conducteur » (titre de
  sa main, précédé d'un 1 cerclé ; il y renvoie plus loin comme au
  « N\textsuperscript{o}~1 »). Un idéal $I$ de $A$ contenu dans le
  conducteur de $A$ dans $A'$ ; la suite exacte qui compare le complexe
  d'Amitsur de $A'/A$ à celui de sa réduction modulo $I$ ; la descente
  effective des modules plats quand $I$ est nilpotent ; le transport des
  relations d'équivalence.
  \item[II.] \emph{Pages 12 à 18}, « Résultats auxiliaires sur les relations
  d'équivalence » (2 cerclé). Les extensions $A \subset A'$ où $A'/A$ est le
  corps résiduel ; un critère pour que $A'/A$ soit de longueur au plus $1$,
  dont la démonstration s'arrête au dernier cas.
  \item[III.] \emph{Pages 19 à 29}, sans titre de sa main. Le lemme~1 et ses
  corollaires, qui comparent $\mathbb{G}_m$ (puis $\mathrm{Gl}_n$) et
  $\mathbb{G}_a$ (puis $\mathbb{M}_n$) à un idéal de carré nul près ; le
  lemme~2 ; le théorème sur un anneau semi-local noethérien et ses
  voisinages infinitésimaux, avec sa démonstration, inachevée.
\end{itemize}
La page 30 est une chemise vierge ; la page 31, d'une autre main, est décrite
à la fin.

Les trois suites ont une méthode commune, qui fait l'unité du dossier : on
compare le complexe d'Amitsur de $A'/A$ à celui de sa réduction modulo un
idéal nilpotent, par une suite exacte courte de complexes dont le noyau est
simple à calculer. En I le noyau est un complexe constant, et c'est l'hypothèse
du conducteur qui le rend tel ; en III il est seulement de carré nul, et la
comparaison devient une comparaison entre un groupe multiplicatif et un groupe
additif. La suite II sert la suite I : ses deux propositions sont démontrées
en se ramenant au N\textsuperscript{o}~1 avec $I = \mathfrak{m}$.

\subsection*{Conventions}

\begin{itemize}
  \item $A'$ désigne toujours la $A$-algèbre qu'on veut descendre ; les pages
  1 à 11 l'appellent $B$\footnote{Il écrit « $B = A'$ » à la page 4, puis
  $A'$ seul à partir de la page 12. Comme les pages 19 à 29 réservent $B$ à
  l'anneau $H^0(A'/A)$, on garde ici cet usage pour tout le dossier ; sur
  les pages 1 à 11, son $B$ est notre $A'$ et son $B_0$ notre $A'_0$.}.
  Pour un idéal $I$ de $A$ : $A_0 = A/I$, $A'_0 = A'/IA'$.
  \item $C^n(A'/A) = A'^{\otimes_A (n+1)} = \bigotimes^{n+1}_A A'$ ; ces
  anneaux forment le \emph{complexe d'Amitsur} $C^{\bullet}(A'/A)$, et
  $H^i(A'/A, M)$ désigne la cohomologie de $C^{\bullet}(A'/A) \otimes_A M$ ;
  $H^i(A'/A) = H^i(A'/A, A)$. Ainsi $H^0(A'/A) = \mathrm{Ker}(A'
  \rightrightarrows A' \otimes_A A')$. On note $B = H^0(A'/A)$ et
  $\widetilde{B}_0 = H^0(A'_0/A_0)$\footnote{Les pages 19 à 22 écrivent
  $C_0$ pour $H^0(A'_0/A_0)$ ; on l'évite, parce que $C$ sert déjà au
  complexe.}.
  \item Pour un faisceau de groupes $G$ ($\mathbb{G}_a$, $\mathbb{G}_m$,
  $\mathrm{Gl}_n$, $\mathbb{M}_n$), $H^1(A'/A, G)$ est la cohomologie
  d'Amitsur à coefficients dans $G$ : pour $G$ non commutatif, l'ensemble
  pointé des classes de cocycles $g \in G(A' \otimes_A A')$, et pour
  $\mathrm{Gl}_n$ l'ensemble des classes de données de descente sur
  $A'^{\,n}$.
  \item Une \emph{donnée de descente} sur un $A'$-module $E'$ relativement à
  $A'/A$ est un isomorphisme $E' \otimes_{A', p_1} (A' \otimes_A A') \simeq
  E' \otimes_{A', p_2} (A' \otimes_A A')$ satisfaisant la condition de
  cocycle ; elle est \emph{effective} s'il existe un $A$-module $E$ et un
  isomorphisme $E \otimes_A A' \simeq E'$ compatible aux données. Le
  morphisme $\mathrm{Spec}\, A' \to \mathrm{Spec}\, A$ est \emph{de
  descente} (pour une classe de modules) si le foncteur $E \mapsto E
  \otimes_A A'$ est pleinement fidèle à valeurs dans les modules munis d'une
  donnée de descente, et \emph{de descente effective} si c'est de plus une
  équivalence.
  \item $\mathfrak{m}$ désigne l'idéal maximal (section II) ou le radical
  (section III) ; l'entier qu'il appelle $m$ dans le théorème de la page 25
  est noté $r$\footnote{La page écrit « Soit $m$ un entier $\geqslant 1$ »
  quelques lignes après « $\mathfrak{m}$ son radical » ; la gothique les
  distingue sur la page, mais pas toujours à la lecture.}.
\end{itemize}

\subsection*{Ce que le dossier annonce sans l'établir}

Le critère d'effectivité de la page 11 s'arrête sur des points de suspension ;
le cas $k[\varepsilon]$ de la page 13 n'est pas vérifié ; la page 18 s'arrête
sur ce qu'« il faut prouver ». Ces trois points se démontrent sans peine et on
le fait, en le disant. Plus sérieux : la démonstration du lemme~1 (page 20)
repose sur un carré qui ne commute pas tel qu'il est dessiné, et ses
corollaires (pages 21 et 22), puis les étapes du théorème qui les utilisent,
héritent de cette lacune. La démonstration du lemme~2 et de sa variante est
en grande partie illisible ; la troisième partie du théorème, c), n'est pas
démontrée ; le lemme de passage à la limite de la page 29 n'est pas écrit.

\pagerange{1}{11}
\section*{I. Réductions via le conducteur (pages 1 à 11)}

\pagerange{1}{2}
\subsection*{Le cadre : un idéal contenu dans le conducteur (pages 1 et 2)}

Soient $A$ un anneau, $I$ un idéal de $A$ et $A'$ une $A$-algèbre. On suppose
que l'application $I \to A'$, $\lambda \mapsto \lambda \cdot 1_{A'}$, est
injective et que son image est un idéal de $A'$ — nécessairement égal à
$IA'$\footnote{La précision « nécessairement égal à $IB$ » est ajoutée en
interligne ; une première rédaction, qui supposait $B$ « contenant $A$ », est
rayée.}. Quand $A \subset A'$, cela revient à dire que $I$ est un idéal de
$A'$ contenu dans $A$, c'est-à-dire que $I$ est contenu dans le
\emph{conducteur} $\mathfrak{f} = \mathrm{Ann}_A(A'/A)$ de $A$ dans $A'$ ; le
carré

\begin{tikzcd}
A \arrow[r] \arrow[d] & A' \arrow[d] \\
A/I \arrow[r] & A'/I
\end{tikzcd}

est alors cartésien. C'est la situation qu'on appelle aujourd'hui un carré de
Milnor, ou un pincement\footnote{Le nom vient de Milnor, \emph{Introduction
to algebraic $K$-theory} (1971), qui recolle les modules projectifs le long
d'un tel carré ; l'étude systématique de la descente dans ce cadre est due à
D.~Ferrand, « Conducteur, descente et pincement » (\emph{Bull. SMF}, 2003).
Les feuillets ne citent ni l'un ni l'autre, et ne nomment le conducteur que
dans leur titre.}. L'exemple type est la normalisation d'une courbe à point
double : $A' = k[t]$, $A = \{ f : f(0) = f(1) \}$, $I = \mathfrak{f} = \{ f :
f(0) = f(1) = 0 \}$.

\emph{Proposition 1.} Sous ces hypothèses, pour tout $n \geqslant 0$,
l'application $\lambda \mapsto \lambda \cdot 1$ de $I$ dans $C^n =
\bigotimes^{n+1}_A A'$ est injective et son image est un idéal de $C^n$.

En effet, pour $\lambda \in I$ on a $\lambda b_0 = \lambda_0 \cdot 1_{A'}$
avec $\lambda_0 \in I$, d'où
\[
  \lambda (b_0 \otimes b_1 \otimes \cdots \otimes b_n)
  = 1 \otimes \lambda_0 (b_1 \otimes \cdots \otimes b_n)
  = 1 \otimes \mu (1 \otimes \cdots \otimes 1)
  = \mu \cdot 1_{C^n}
\]
avec $\mu \in I$, par récurrence sur $n$\footnote{La deuxième ligne du calcul
de la page 1 est surchargée et sa lecture incertaine ; le calcul donné ici est
celui que les lignes lisibles imposent. Un passage qui annonçait déjà la suite
exacte du corollaire~1 est rayé.}. L'injectivité vient de ce que le composé
$I \to C^n \to A'$ avec la multiplication est l'application donnée, injective
par hypothèse\footnote{La phrase de la page est en partie illisible ; seul
« induit une injection … par hyp. » se lit.}.

Comme $I \cdot 1_{C^n} = IC^n$ et $C^n / IC^n = \bigotimes^{n+1}_{A_0}
A'_0$, on obtient :

\emph{Corollaire 1.} Pour tout ensemble fini $E$, la suite
\[
  0 \to I \to \bigotimes^E_A A' \to \bigotimes^E_{A_0} A'_0 \to 0
\]
est exacte, et fonctorielle en $E$, l'application induite sur $I$ étant
l'identité. On a donc une suite exacte de $A$-modules
\emph{cosimpliciaux}\footnote{La page dit « simpliciaux ». Le foncteur $E
\mapsto \bigotimes^E_A A'$ est covariant en l'ensemble fini $E$, donc définit
un objet cosimplicial : c'est le complexe d'Amitsur. Le mot « simplicial »
pour cette variance n'était pas fixé à l'époque ; on suit l'usage actuel
dans tout le dossier.}
\[
  0 \to I_{\mathrm{cst}} \to C^{\bullet}(A'/A) \to C^{\bullet}(A'_0/A_0) \to 0 ,
\]
où $I_{\mathrm{cst}}$ est le module cosimplicial constant de valeur
$I$\footnote{La page note $C_I$ ce module (« foncteur $E \mapsto I$ »), et
le complexe d'Amitsur d'une lettre capitale grasse qui tient du $C$ et du
$K$ : $C(B/A)$ à la page 2, $K(B/A)$ à la page 3. Un lemme~1 sur deux
algèbres $B$, $C$, encadré et rayé en tête de la page 2, n'est pas celui des
pages 19 et 20.}.

\pagerange{3}{4}
\subsection*{Ce que la réduction ne change pas (pages 3 et 4)}

Le complexe normalisé de $I_{\mathrm{cst}}$ est $I$ en degré $0$ et nul
ensuite ; sa cohomologie est $I$ en degré $0$ et $0$ en degrés $> 0$. La
suite exacte longue donne donc :

\emph{Corollaire 2.} On a $H^i(A'/A) \simeq H^i(A'_0/A_0)$ pour $i \neq 0$,
et une suite exacte $0 \to I \to H^0(A'/A) \to H^0(A'_0/A_0) \to 0$. Par
le lemme du serpent appliqué à $0 \to I \to A \to A_0 \to 0$, on a aussi
$\mathrm{Coker}(A \to H^0(A'/A)) \simeq \mathrm{Coker}(A_0 \to
H^0(A'_0/A_0))$, et de même pour les noyaux\footnote{L'isomorphisme des
conoyaux est écrit en travers de la marge gauche ; celui des noyaux n'est pas
écrit, mais c'est lui qu'utilise le corollaire~3 b).}.

\emph{Corollaire 3.} a) Pour $i \neq 0$, $H^i(A'/A) = 0$ si et seulement si
$H^i(A'_0/A_0) = 0$. b) $A \to H^0(A'/A)$ est injectif (resp. surjectif) si
et seulement si $A_0 \to H^0(A'_0/A_0)$ l'est\footnote{« Cor.~3 » et les
lettres a), b) sont ajoutés dans la marge.}.

\emph{Corollaire 4.} Si $M$ est un $A$-module plat, on a une suite exacte de
modules cosimpliciaux $0 \to (M \otimes_A I)_{\mathrm{cst}} \to
C^{\bullet}(A'/A) \otimes_A M \to C^{\bullet}(A'_0/A_0) \otimes_{A_0} M_0
\to 0$, d'où $H^i(A'/A, M) \simeq H^i(A'_0/A_0, M_0)$ pour $i \neq 0$ et une
suite exacte $0 \to M \otimes_A I \to H^0(A'/A, M) \to H^0(A'_0/A_0, M_0)
\to 0$.

\emph{Corollaire 5.} Si $A'_0$ est fidèlement plat sur $A_0$, alors pour
tout $A$-module plat $M$ on a $H^i(A'/A, M) = 0$ pour $i \neq 0$ et $M
\xrightarrow{\sim} H^0(A'/A, M)$.

C'est la descente fidèlement plate pour $A'_0/A_0$, remontée par le
corollaire~4 et le lemme des cinq. Le résultat équivaut au cas $M = A$ —
c'est-à-dire à l'exactitude du complexe augmenté $0 \to A \to A' \to A'
\otimes_A A' \to \cdots$ —, puisqu'un tel complexe reste exact après
tensorisation par un module plat\footnote{La page écrit « $H^0(B, A)
\xleftarrow{\sim} A$ » dans la seconde accolade ; il faut lire $H^0(B/A)$.}.

\pagerange{4}{9}
\subsection*{Descente effective le long d'un conducteur nilpotent (pages 4 à 9)}

\emph{Proposition 2.} Soient $A$, $I$, $A'$ comme ci-dessus, avec $I$
\emph{nilpotent}, et tels que $A \xrightarrow{\sim} H^0(A'/A)$ et $H^1(A'/A)
= 0$\footnote{« et $I$ nilpotent » est ajouté dans la marge, le dernier mot
souligné deux fois ; la parenthèse des conditions sur $H^0$ et $H^1$ est en
interligne. La note de bas de page de sa main dit que ces conditions
« signifient que $\mathrm{Spec}\, B \to \mathrm{Spec}\, A$ est un morphisme
de descente pour la cat.\ fib.\ des modules plats », et qu'elles équivalent
aux mêmes conditions pour $A'_0/A_0$ (corollaire~3). La seconde affirmation
est exacte. La première ne l'est qu'à moitié : être un morphisme de descente
pour les modules plats équivaut à la seule condition $A \simeq H^0(A'/A)$,
car pour $N$ plat on a $H^0(A'/A, N) = N \otimes_A H^0(A'/A)$ et
$\mathrm{Hom}_A(M, -)$ commute aux noyaux. La condition $H^1(A'/A) = 0$ est
de plus ; la démonstration qui suit ne s'en sert pas.}. Soit $E'$ un
$A'$-module plat muni d'une donnée de descente relativement à $A'/A$, et
supposons effective la donnée qu'elle induit sur $E'_0 = E'/IE'$
relativement à $A'_0/A_0$ : il existe un $A_0$-module plat $E_0$ et un
isomorphisme compatible aux données $E_0 \otimes_{A_0} A'_0 \simeq E'_0$.
Alors la donnée sur $E'$ est effective : si $E = H^0_D(A'/A, E')$ est le
sous-$A$-module des invariants de $E'$, $E$ est plat sur $A$ et $E \otimes_A
A' \to E'$ est un isomorphisme.

La page dessine les deux étages en perspective : en bas les complexes
d'Amitsur de $A'/A$ et de $A'_0/A_0$, en haut les complexes de descente
$C^{\bullet}_D(A'/A, E')$, de termes $E^{(n)} = E' \otimes_{A'} A^{(n)}$ où
$A^{(n)} = \bigotimes^{n+1}_A A'$, et leur réduction, avec $E_0$ en pointillé
au-dessus de $A_0$\footnote{Le dessin est rendu dans la transcription ; ses
traits doubles et triples figurent les deux et trois flèches de faces.}.

\emph{Démonstration, cas $I^2 = 0$.} On a une suite exacte de modules
cosimpliciaux
\[
  0 \to I\, C^{\bullet}_D(A'/A, E') \to C^{\bullet}_D(A'/A, E') \to
  C^{\bullet}_D(A'_0/A_0, E'_0) \to 0 ,
\]
qui n'utilise pas l'effectivité\footnote{La page écrit « sans hypothèse
d'effectivité pour $E'_0$ rel.\ à $A'_0/A'$ » ; il faut lire $A'_0/A_0$.}.
Comme $E^{(n)}$ est plat sur $A^{(n)}$ et $I^2 = 0$, on a $I E^{(n)} \simeq
E^{(n)}_0 \otimes_{A^{(n)}_0} IA^{(n)}$, et $IA^{(n)} \simeq I$ par la
proposition~1. L'effectivité modulo $I$ donne $E^{(n)}_0 \simeq E_0
\otimes_{A_0} A^{(n)}_0$, d'où $IE^{(n)} \simeq E_0 \otimes_{A_0} I$, et ces
isomorphismes sont compatibles aux opérations cosimpliciales : le noyau est le
module constant $(E_0 \otimes_{A_0} I)_{\mathrm{cst}}$. La suite exacte de
cohomologie donne donc
\[
  0 \to E_0 \otimes_{A_0} I \to E \to H^0_D(A'_0/A_0, E'_0) \to 0 ,
\]
et $H^0_D(A'_0/A_0, E'_0) = E_0 \otimes_{A_0} H^0(A'_0/A_0) = E_0$ puisque
$E_0$ est plat et $H^0(A'_0/A_0) = A_0$. Ainsi $E \to E_0$ est surjectif,
son noyau est l'image de $E_0 \otimes_{A_0} I$, qui est $IE$, et :

a) $E/IE \xrightarrow{\sim} E_0$ ;

b) $E/IE \otimes_{A_0} I = E \otimes_A I \to IE$ est un isomorphisme.

Comme $E/IE$ est plat sur $A/I$ et $I$ nilpotent, b) dit que
$\mathrm{Tor}^A_1(A/I, E) = 0$, et $E$ est plat sur $A$ par le critère local
de platitude, sous sa forme pour un idéal nilpotent, qui ne demande aucune
hypothèse noethérienne\footnote{La page conclut « donc grâce à a) et $E_0$
plat sur $A_0$, $E$ est \emph{plat} sur $A$ », sans nommer le critère.}.
Enfin $E \otimes_A A' \to E'$ est un homomorphisme de $A'$-modules plats qui
devient un isomorphisme modulo $I$ ; il est surjectif par Nakayama ($IA'$
est nilpotent), et son noyau $K$ vérifie $K/IK = 0$ parce que $E'$ est plat,
donc $K = 0$\footnote{Ces deux lignes de la page 6 sont presque entièrement
illisibles ; une première rédaction, rayée, invoquait déjà « Nakayama sur
$A'$ ». L'argument donné ici est le nôtre, et c'est celui que les mots lisibles
(« $IA'$ », « donc est un isomorphisme ») laissent attendre.}.

\emph{Cas général.} Si $I^n = 0$ avec $n > 2$, on raisonne par récurrence sur
$n$, au moyen du lemme suivant, que la page dit « trivial » :

\emph{Lemme.} Si $I$ et $J$ sont deux idéaux de $A$ dont les images
$\lambda \mapsto \lambda \cdot 1$ dans $A'$ sont injectives et sont des idéaux
de $A'$, il en est de même de $IJ$ ; en particulier des puissances $I^{\nu}$,
$\nu \geqslant 1$.

(Pour $b \in A'$, $\lambda \in I$, $\mu \in J$ : $b \lambda\mu = \lambda (\mu
b) = \lambda \mu'$ avec $\mu' \in J$.) On applique alors le cas de carré nul
à l'idéal $J = I^{n-1}$, de carré nul, au-dessus de $A/J$, où le résultat est
acquis par récurrence pour l'idéal $I/J$, nilpotent d'ordre $n-1$\footnote{La
page ne détaille pas la récurrence ; on en écrit le découpage, qui est le
seul que le lemme rende possible.}.

\emph{Corollaire 1.} Soient $A \subset A'$ et $I$ un idéal nilpotent de $A$
tel que $IA' = I$ (c'est-à-dire contenu dans le conducteur). Si
$\mathrm{Spec}\, A'_0 \to \mathrm{Spec}\, A_0$ est un morphisme de descente
effective pour les modules plats — par exemple si $A'_0$ est fidèlement plat
sur $A_0$ —, il en est de même de $\mathrm{Spec}\, A' \to \mathrm{Spec}\,
A$\footnote{« morphisme de » et « p.ex.\ $B_0/A_0$ fid.\ plat » sont en
interligne, au-dessus d'un « fid.\ plat » rayé : l'hypothèse a été
affaiblie en cours d'écriture.}.

\emph{Corollaire 2.} Soient $X$ un préschéma, $\mathcal{I}$ un idéal
quasi-cohérent nilpotent de $\mathcal{O}_X$, et $X' =
\mathrm{Spec}(\mathcal{A}') \to X$ affine, tel que $\mathcal{I} \to
\mathcal{A}'$ soit un isomorphisme sur un idéal de $\mathcal{A}'$ et que
$X'_0 = \mathrm{Spec}(\mathcal{A}'/\mathcal{I}\mathcal{A}')$ soit
fidèlement plat sur $X_0 = V(\mathcal{I})$. Alors $X' \to X$ est un
morphisme de descente effective pour les modules quasi-cohérents plats.

\emph{Corollaire 3.} Sous les mêmes conditions, pour tout morphisme
\emph{plat} $Y \to X$, la projection $X' \times_X Y \to Y$ est encore un
morphisme de descente effective pour les modules quasi-cohérents plats.

En effet les hypothèses du corollaire~2 sont stables par changement de base
plat\footnote{« plat » est ajouté au crayon dans une bulle ; la page appelle
cette propriété « universelle ». Elle ne l'est que pour les changements de base
plats, ce qui est ce que dit l'énoncé.}. La page ajoute que cela s'applique
« lorsque $Y$ est un $A$-préschéma quelconque » ; il faut lire : un
$A$-préschéma \emph{plat}. Pour $Y$ quelconque l'énoncé est faux, parce que
le carré cartésien du conducteur ne le reste pas par un changement de base
non plat\footnote{Contre-exemple (le nôtre) : $A' = k[t]/(t^2) \times
k[t]/(t^2)$, $u = (t, 0)$, $v = (0, t)$, $A = k \oplus ku \oplus kv$ et $I
= (u, v)$, de carré nul, idéal de $A'$ ; $A'_0 = k \times k$ est
fidèlement plat sur $A_0 = k$. Pour $R = A/(u - v)$, on trouve $R \otimes_A
A' = A'/(u - v)A' = k \times k$, et $R = k[u]/(u^2) \to k \times k$ envoie
$u$ sur $0$ : il n'est pas injectif, et $\mathrm{Spec}(R \otimes_A A') \to
\mathrm{Spec}\, R$ n'est même pas un morphisme de descente.}.

\pagerange{9}{11}
\subsection*{Idéaux d'augmentation et relations d'équivalence (pages 9 à 11)}

Revenons au diagramme du corollaire~1 de la proposition~1, avec pour $F$ un
point : la multiplication $\bigotimes^n_A A' \to A'$ et sa réduction
s'inscrivent dans un morphisme de suites exactes $0 \to I \to
\bigotimes^n_A A' \to \bigotimes^n_{A_0} A'_0 \to 0$ vers $0 \to I \to A'
\to A'_0 \to 0$, l'identité sur $I$. Le lemme du serpent donne :

\emph{Proposition 3.} Soit $\Delta^{(n)}_{A'/A} = \mathrm{Ker}\bigl(
\bigotimes^n_A A' \to A'\bigr)$ l'idéal d'augmentation. La réduction modulo
$I$ induit un isomorphisme $\Delta^{(n)}_{A'/A} \xrightarrow{\sim}
\Delta^{(n)}_{A'_0/A_0}$\footnote{Un premier « Proposition 3 », en tête du
paragraphe, est rayé et l'énoncé récrit après le diagramme ; « $A$, $I$, $B$
comme dans Prop.~1 » est en interligne.}.

\emph{Corollaire 1.} Notons $p$ la réduction. Les applications $J \mapsto
p(J)$ et $J^0 \mapsto p^{-1}(J^0) \cap \Delta^{(n)}_{A'/A}$ sont des
bijections réciproques entre idéaux de $\bigotimes^n_A A'$ contenus dans
l'idéal d'augmentation et idéaux de $\bigotimes^n_{A_0} A'_0$ contenus dans
le sien ; pour deux idéaux qui se correspondent, $J \to J^0$ est bijectif.
La correspondance respecte l'inclusion, les sommes et les intersections, et
commute aux opérations cosimpliciales (pour $n$ variable)\footnote{La lecture
de « $\cap$ » est incertaine ; l'énoncé est vrai avec lui.}.

Une \emph{relation d'équivalence} de $\mathrm{Spec}\, A'$ au-dessus de
$\mathrm{Spec}\, A$ est un sous-schéma fermé de $\mathrm{Spec}(A' \otimes_A
A')$, défini par un idéal $J \subset \Delta^{(2)}$ (réflexivité), stable par
l'échange des facteurs (symétrie) et tel que $p_{13}(J) \subset p_{12}(J) +
p_{23}(J)$ dans $A'^{\otimes 3}$ (transitivité) : un \emph{idéal
d'équivalence}. Ces conditions s'expriment par sommes, inclusions et
opérations cosimpliciales ; d'où :

\emph{Corollaire 2.} (a) La correspondance du corollaire~1 (pour $n = 2$) est
une bijection entre relations d'équivalence de $\mathrm{Spec}\, A'$ sur
$\mathrm{Spec}\, A$ et de $\mathrm{Spec}\, A'_0$ sur $\mathrm{Spec}\, A_0$.
(b) Elle respecte l'ordre ; la plus fine (la diagonale, $J = \Delta$) et la
moins fine ($J = 0$, le produit fibré tout entier) se correspondent. (c) Si
$J$ et $J^0$ se correspondent, et si $A'^{J} = \{ b : p_1(b) \equiv p_2(b)
\bmod J \}$ désigne l'anneau des invariants, alors $A'^{J} \to A'^{J^0}_0$
est surjectif de noyau $I$ : $A'^{J^0}_0 \simeq A'^{J}/I$. En effet $\delta(b)
= p_2(b) - p_1(b)$ appartient à $\Delta$, et $\Delta/J \simeq \Delta_0/J^0$,
de sorte que $\delta(b) \in J$ si et seulement si $\delta_0(b_0) \in J^0$.
Enfin $J$ est effectif si et seulement si $J^0$ l'est\footnote{L'alinéa (c)
est entre crochets sur la page et s'arrête sur « il f.\ et s.\ que $J^0$ le
soit … ] » ; la marge porte « expliciter ». La justification est la nôtre : si
« effectif » veut dire que $J$ est l'idéal engendré par les $\delta(c)$, $c
\in A'^{J}$ — c'est-à-dire que la relation est le graphe de $\mathrm{Spec}\,
A' \to \mathrm{Spec}\, A'^{J}$ —, cet idéal correspond, par la surjection
$A'^{J} \to A'^{J^0}_0$, à l'idéal engendré par les $\delta_0(c_0)$.}.

\pagerange{12}{18}
\section*{II. Résultats auxiliaires sur les relations d'équivalence (pages 12 à 18)}

\pagerange{12}{13}
\subsection*{Les extensions où $A'/A$ est le corps résiduel (pages 12 et 13)}

\emph{Proposition 1.} Soient $A$ un anneau, $\mathfrak{m}$ un idéal maximal,
$k = A/\mathfrak{m}$, et $A \subset A'$ une $A$-algèbre telle que $A'/A
\simeq k$ comme $A$-modules. Alors :

(i) il y a exactement deux relations d'équivalence de $\mathrm{Spec}\, A'$
sur $\mathrm{Spec}\, A$, la plus fine et la moins fine ;

(ii) la suite $0 \to A \to A' \rightrightarrows A' \otimes_A A'$ est exacte,
et plus précisément $H^0(A'/A) = A$, $H^i(A'/A) = 0$ pour $i \neq 0$ ;

(iii) si de plus $\mathfrak{m}$ est nilpotent, $\mathrm{Spec}\, A' \to
\mathrm{Spec}\, A$ est un morphisme de descente effective pour les modules
quasi-cohérents plats, et le reste après tout changement de base plat.

Comme $\mathfrak{m}$ annule $A'/A$, on a $\mathfrak{m}A' \subset A$ ; c'est un
idéal de $A$ contenant $\mathfrak{m}$, distinct de $A$ (sinon $A' =
\mathfrak{m}A' \subset A$), donc $\mathfrak{m}A' = \mathfrak{m}$ : on est dans
la situation du N\textsuperscript{o}~1 avec $I = \mathfrak{m}$. De plus
$A'_0 = A'/\mathfrak{m}$ est de dimension $2$ sur $k$, donc fidèlement plat :
(ii) résulte du corollaire~5 de la proposition~1, et (iii) du corollaire~1 de
la proposition~2 et du corollaire~3\footnote{La page énonce (iii) sans
supposer $\mathfrak{m}$ nilpotent, et pour « toute extension de la base $X \to
\mathrm{Spec}(A)$ ». Or la proposition~2 du N\textsuperscript{o}~1 exige un
idéal nilpotent, et son corollaire~3 un changement de base plat ; c'est donc
sous ces deux restrictions que (iii) est démontré. Pour $A$ quelconque —
la courbe à point double, par exemple, où $A'/A$ est bien le corps résiduel
du point double — l'énoncé relève de la théorie du pincement, que le dossier
ne développe pas, et on ne le vérifie pas ici. Dans la suite (proposition~2),
$A$ est local artinien et la restriction ne coûte rien.}.

Pour (i), la proposition~3 et son corollaire~2 ramènent au cas $A = k$, où
$A'$ est une $k$-algèbre de rang $2$ ; on peut supposer $k$ algébriquement
clos\footnote{En interligne, « On peut supposer $k$ alg.\ clos ». C'est
légitime : la plus fine et la moins fine des relations sont toujours
distinctes, et une extension du corps de base fait des relations
d'équivalence sur $k$ une partie de celles sur $\bar{k}$.}. Deux cas :

(1) $A'$ réduit, donc $A' \simeq k \times k$ : la question est ensembliste,
et un ensemble à deux éléments a deux relations d'équivalence.

(2) $A'$ non réduit, donc $A' \simeq k[\varepsilon]$, $\varepsilon^2 = 0$.
La page s'arrête là ; la vérification est courte\footnote{Elle est la nôtre.
Après « $\varepsilon^2 = 0$ » la page ne porte que deux mots illisibles entre
parenthèses.}. Dans $A' \otimes_k A' = k[\varepsilon_1, \varepsilon_2]$,
l'idéal d'augmentation est $\Delta = k(\varepsilon_1 - \varepsilon_2) \oplus
k\varepsilon_1\varepsilon_2$, et ses idéaux sont $0$, $k\varepsilon_1
\varepsilon_2$ et $\Delta$ (tout élément de $\Delta$ hors de $k\varepsilon_1
\varepsilon_2$ engendre $\Delta$). L'idéal $k\varepsilon_1\varepsilon_2$
n'est pas transitif : $\varepsilon_1\varepsilon_3$ n'appartient pas à
l'idéal engendré par $\varepsilon_1\varepsilon_2$ et $\varepsilon_2
\varepsilon_3$ dans $k[\varepsilon_1, \varepsilon_2, \varepsilon_3]$. Il
reste $0$ et $\Delta$.

\pagerange{14}{18}
\subsection*{Un critère de colongueur au plus un (pages 14 à 18)}

\emph{Proposition 2.} Soient $A$ un anneau local artinien d'idéal maximal
$\mathfrak{m}$, $A'$ une $A$-algèbre finie dont les extensions résiduelles
sont triviales, $A'' = (A' \otimes_A A')/J$ avec $J$ un idéal d'équivalence,
et $\Delta = \mathrm{Ker}(A'' \to A')$. On suppose que la suite $A \to A'
\rightrightarrows A''$ est exacte, et que pour tout $x' \in A' \smallsetminus
A$, l'élément $\delta(x') = p_2(x') - p_1(x')$ engendre l'idéal $\Delta$.
Alors $A'/A$ est de longueur au plus $1$ ; si elle vaut $1$, on a $A''
\simeq A' \otimes_A A'$\footnote{L'hypothèse sur les extensions résiduelles
est ajoutée dans la marge, appelée devant « Sous ces conditions » ; on l'intègre
à l'énoncé, puisque la démonstration s'en sert dans les deux cas. Une seconde
note marginale, de lecture douteuse, dit que sans elle on aurait encore
$\mathfrak{m}A' = \mathfrak{m}$ et qu'on se ramènerait à $A = k$ ; la page ne
la développe pas. La dernière assertion vient de la proposition~1 (i) : $J =
\Delta$ rendrait $A'' = A'$ et contredirait l'exactitude.}.

Notons $\mathfrak{m}' = \mathfrak{r}(A')$, le radical, formé d'éléments
nilpotents puisque $A'$ est artinien. Deux cas.

\emph{1\textsuperscript{o} $A'$ n'est pas radiciel sur $A$}, c'est-à-dire
(les extensions résiduelles étant triviales) a au moins deux idéaux maximaux.
Alors $\mathrm{Ker}(A' \otimes_A A' \to A')$ n'est pas nilpotent. Comme $J$
est nilpotent, $\Delta$ ne l'est pas non plus\footnote{« Comme $J$ est
nilpotent » est souligné sur la page, sans justification. En voici une, la
nôtre : chaque composante $\mathrm{Spec}(A'_i \otimes_A A'_j)$ du produit
est réduite à un point (extensions résiduelles triviales) ; la relation
définie par $J$ induit une relation d'équivalence sur l'ensemble des
composantes locales de $A'$, et l'idempotent d'une classe est invariant ;
l'exactitude force donc une seule classe, la relation rencontre chaque
composante, donc la contient ensemblistement, et $J$ est nilpotent.}. Si
$x' \in \mathfrak{m}'$, $\delta(x')$ est nilpotent, donc n'engendre pas
$\Delta$ ; par hypothèse $x' \in A$. Ainsi $\mathfrak{m}' \subset A$, donc
$\mathfrak{m}' \subset \mathfrak{m}$, et comme $\mathfrak{m} \subset
\mathfrak{m}A' \subset \mathfrak{m}'$,
\[
  \mathfrak{m}' = \mathfrak{m} = \mathfrak{m}A' .
\]
On est dans la situation du N\textsuperscript{o}~1 avec $I = \mathfrak{m}$ et
$A_0 = k$. Par la proposition~3 (corollaire~2 (c)), la suite réduite $k \to
A'_0 \rightrightarrows A''_0$ est encore exacte ; $A'_0 = A'/\mathfrak{m}'$ est
réduit ; et comme $\Delta \to \Delta_0 = \mathrm{Ker}(A''_0 \to A'_0) \simeq
\Delta/(\Delta \cap \mathfrak{m}A'')$ est surjectif, $\delta_0(x'_0)$ engendre
$\Delta_0$ pour tout $x'_0 \in A'_0 \smallsetminus k$\footnote{La page écrit
« $\delta_0(x'_0) \in A'_0$ » ; l'élément est dans $\Delta_0 \subset
A''_0$.}. Or $A'_0 \simeq k^n$ avec $n \geqslant 2$, l'exactitude impose
$A''_0 = k^{n^2}$ (une seule classe), et $\Delta_0$, l'idéal des fonctions
nulles sur la diagonale, est de dimension $n^2 - n$. Pour $x'_0 = (1, 0,
\ldots, 0)$, $\delta_0(x'_0)$ est non nul exactement aux $2(n-1)$ couples
$(i, j)$ dont une seule coordonnée vaut $1$ ; il engendre $\Delta_0$ si et
seulement si $2(n-1) = n^2 - n$, c'est-à-dire $n = 2$. Alors $A'/A \simeq
A'_0/k \simeq k$, de longueur $1$\footnote{La page conclut « on trouve $n =
2$, OK » ; le décompte est le nôtre.}.

\emph{2\textsuperscript{o} $A'$ est radiciel sur $A$} : $\mathrm{Ker}(A'
\otimes_A A' \to A')$ est nilpotent, donc $\Delta$ l'est. Posons $\Omega =
\Delta/\Delta^2$.

a) Si $\Omega = 0$, alors $\Delta = \Delta^2$, donc $\Delta = 0$ ; $A'' =
A'$, et l'exactitude donne $A = A'$ : la longueur est nulle.

b) Si $\Omega \neq 0$. L'application $d : A' \to \Omega$, $d(x') =
\delta(x') \bmod \Delta^2$, est une $A$-dérivation — l'image de la
différentielle universelle par l'application canonique $\Omega^1_{A'/A} \to
\Omega$ —, donc $d(\mathfrak{m}'^2) \subset \mathfrak{m}'\Omega$. Par
hypothèse, $d(x')$ engendre le $A'$-module $\Omega$ pour $x' \notin A$ ; par
Nakayama, un élément de $\mathfrak{m}'\Omega$ n'engendre pas $\Omega$. Donc
\[
  \mathfrak{m}'^2 \subset A , \quad \text{d'où} \quad \mathfrak{m}'^2 \subset
  \mathfrak{m} .
\]
Comme $A'$ est local à extension résiduelle triviale, $A' = A +
\mathfrak{m}'$, d'où $\mathfrak{m}A' \subset \mathfrak{m} + \mathfrak{m}'^2
\subset \mathfrak{m}$, soit $\mathfrak{m}A' = \mathfrak{m}$ : on est de
nouveau dans la situation du N\textsuperscript{o}~1, et l'on se ramène comme
au 1\textsuperscript{o} au cas où $A = k$ est un corps\footnote{En marge, en
oblique : « vérifier les préliminaires » (lecture en partie incertaine).}.
Alors $A'$ est local d'idéal maximal $V$ de carré nul,
\[
  A' \simeq D_k(V) = k \oplus V , \quad V^2 = 0 ,
\]
l'algèbre des nombres duaux sur un espace $V$ de dimension finie, non nul
puisque $\Delta \neq 0$.

Il faut prouver que $\dim V = 1$ ; la page s'arrête après en avoir préparé
le calcul\footnote{La page 18 s'arrête, à mi-hauteur, sur la description de
l'idéal engendré par $\delta(v)$ ; elle écrit « les multiples de
$\delta(w)$ », qu'il faut lire $\delta(v)$. La fin de la démonstration est la
nôtre. Elle n'utilise de $J$ que le fait d'être un idéal contenu dans
l'idéal d'augmentation ; la feuille d'une autre main de la page 31 traite des
idéaux d'équivalence de ce même $k \oplus V$.}. On a $A' \otimes_k A' = k
\oplus (V \otimes 1) \oplus (1 \otimes V) \oplus (V \otimes V)$, l'idéal
d'augmentation est $\Delta^{(2)} = \delta(V) \oplus (V \otimes V)$, et pour
$v \in V$ l'idéal engendré par $\delta(v) = 1 \otimes v - v \otimes 1$ est
$k\,\delta(v) + V \otimes v + v \otimes V$. Soit $W \subset V$ l'image de $J$
dans $\Delta^{(2)}/(V \otimes V) \simeq V$. L'hypothèse dit que pour tout $v
\neq 0$, $(\delta(v)) + J = \Delta^{(2)}$, donc $V = kv + W$. Si $W$
contenait un $w \neq 0$, on aurait $V = W$ ; alors pour tout $w \in V$, $J$
contiendrait un élément $\delta(w) + t$ avec $t \in V \otimes V$, et en le
multipliant par $u \otimes 1$ et $1 \otimes u$, tous les $u \otimes w$ et $w
\otimes u$ : $V \otimes V \subset J$, puis $\delta(V) \subset J$, ce qui
contredit l'exactitude ($\delta(v) \in J$ voudrait dire $v \in k$). Donc $W
= 0$, et $V = kv$ pour tout $v \neq 0$ : $\dim V = 1$, et $A'/A$ est de
longueur $1$.

\pagerange{19}{29}
\section*{III. Descente le long des voisinages infinitésimaux (pages 19 à 29)}

\pagerange{19}{20}
\subsection*{Le lemme 1 : le groupe multiplicatif et le groupe additif (pages 19 et 20)}

Le cadre change : $I$ n'est plus supposé contenu dans le conducteur, mais de
carré nul.

\emph{Lemme 1.} Soient $A$ un anneau, $A'$ une $A$-algèbre, $I$ un idéal de
carré nul de $A$, $B = H^0(A'/A)$, $\widetilde{B}_0 = H^0(A'_0/A_0)$, et
$\overline{B}$ l'image de $B$ dans $\widetilde{B}_0$. On suppose
\[
  \widetilde{B}_0 = \widetilde{B}_0^{\,*} + \overline{B} .
\]
Le manuscrit affirme qu'il existe alors un isomorphisme canonique
\[
  (*) \qquad \mathrm{Ker}\bigl(H^1(A'/A, \mathbb{G}_m) \to H^1(A'_0/A_0,
  \mathbb{G}_m)\bigr) \simeq \mathrm{Ker}\bigl(H^1(A'/A, \mathbb{G}_a) \to
  H^1(A'_0/A_0, \mathbb{G}_a)\bigr) .
\]
C'est l'énoncé de la page 19, sur un « Proposition » rayé\footnote{« Lemme 1 » est écrit au-dessus de « Proposition », rayé ; deux
rédactions des hypothèses, l'une par $H^1(A'_0/A_0) = e$ et $A_0 \simeq
H^0(A'_0/A_0)$, l'autre par $H^1(A'_0/A_0, G_a) = 0$ et $H^1(A'_0/A_0,
G_m)$, sont rayées.}

\emph{N.B.} La condition est remplie si $A$ est noethérien semi-local et $A'$
fini sur $A$ : alors $\mathrm{Spec}\, \widetilde{B}_0 \to \mathrm{Spec}\,
\overline{B}$ est radiciel et surjectif, $\widetilde{B}_0$ et $\overline{B}$
sont semi-locaux, $\mathrm{Max}\, \widetilde{B}_0 \to \mathrm{Max}\,
\overline{B}$ est bijectif, et par le lemme chinois on trouve, pour $c \in
\widetilde{B}_0$, un $b \in \overline{B}$ tel que $c - b$ ne soit dans aucun
idéal maximal\footnote{La page dit « d'où facilement l'assertion faite » ; le
recours au lemme chinois est le nôtre, et il n'utilise que l'injectivité sur
les idéaux maximaux. Que $\mathrm{Spec}\, \widetilde{B}_0 \to
\mathrm{Spec}\, \overline{B}$ soit radiciel est affirmé sur la page, non
démontré, et on ne le reconstruit pas.}.

\emph{Ce que la démonstration établit.} Le complexe d'Amitsur et sa
réduction donnent deux suites exactes de modules cosimpliciaux
\[
  0 \to I\,C^{\bullet}(A'/A) \to C^{\bullet}(A'/A) \to C^{\bullet}(A'_0/A_0)
  \to 0 ,
\]
\[
  0 \to I\,C^{\bullet}(A'/A) \to C^{\bullet}(A'/A)^{*} \to
  C^{\bullet}(A'_0/A_0)^{*} \to 0 ,
\]
la seconde par $x \mapsto 1 + x$, qui est un morphisme de groupes parce que
$I^2 = 0$ ; la surjectivité vient de ce qu'une unité se relève modulo un
idéal nilpotent. Posons $H = H^1(I\,C^{\bullet}(A'/A))$. Les suites exactes
de cohomologie identifient les deux membres de $(*)$ respectivement à
\[
  \mathrm{Coker}\bigl(\partial_m : \widetilde{B}_0^{\,*} \to H\bigr)
  \quad \text{et} \quad
  \mathrm{Coker}\bigl(\partial_a : \widetilde{B}_0 \to H\bigr) ,
\]
où $\partial_a(c)$ est la classe de $p_2(\tilde{c}) - p_1(\tilde{c})$ pour un
relèvement $\tilde{c} \in A'$, et $\partial_a$ s'annule sur $\overline{B}$.
Jusqu'ici tout est exact.

\emph{La lacune.} Le manuscrit conclut par un carré commutatif où
$\widetilde{B}_0^{\,*} \to H$ et $\widetilde{B}_0 \to H$ sont reliés par
l'inclusion $\widetilde{B}_0^{\,*} \subset \widetilde{B}_0$ et l'identité de
$H$ ; avec $\widetilde{B}_0 = \widetilde{B}_0^{\,*} + \overline{B}$, les deux
conoyaux seraient égaux. Mais ce carré ne commute pas : pour $u \in
\widetilde{B}_0^{\,*}$ relevé en $\tilde{u}$, le cobord multiplicatif est
$p_2(\tilde{u})\, p_1(\tilde{u})^{-1} = 1 + (p_2(\tilde{u}) -
p_1(\tilde{u}))\, p_1(\tilde{u})^{-1}$, de sorte que
\[
  \partial_m(u) = u^{-1} \cdot \partial_a(u) ,
\]
où $\widetilde{B}_0$ opère sur $H$ par cup-produit en degré $0$ : c'est la
dérivée logarithmique, non la dérivée. Les deux images $\partial_a(
\widetilde{B}_0)$ et $\partial_m(\widetilde{B}_0^{\,*})$ n'ont pas de raison
générale de coïncider\footnote{La page dit des inclusions $i$, $j$ qu'elles
« ne sont pas multiplicatives, mais cela est sans importance ». Le défaut
n'est pas là, mais dans l'identité verticale de droite. Que l'argument formel
ne puisse pas suffire se voit sur un modèle : une dérivation $D$ de $K =
k(u)$, $u^p = a \in k$ en caractéristique $p$, dans $\Omega^1_{K/k}$, nulle
sur $k$, où $K = K^* + k$ ; $u^{-1}du$ n'est pas dans $D(K)$. On ne sait pas
si un tel modèle provient d'un complexe d'Amitsur, et donc pas si le lemme
est faux ; seulement que sa démonstration ne l'établit pas.}. Le lemme~1 est
donc, dans ce dossier, un énoncé annoncé et non démontré ; ce qui est acquis
est la description des deux noyaux comme conoyaux de $\partial_m$ et
$\partial_a$, reliés par $\partial_m(u) = u^{-1}\partial_a(u)$.

\emph{Corollaire 1} (annoncé). Si $H^1(A'_0/A_0, \mathbb{G}_a) = 0$ et
$H^1(A'_0/A_0, \mathbb{G}_m) = 0$, on aurait un isomorphisme canonique
$H^1(A'/A, \mathbb{G}_m) \simeq H^1(A'/A, \mathbb{G}_a)$, de sorte que
$H^1(A'/A, \mathbb{G}_m) = 0$ et $H^1(A'/A, \mathbb{G}_a) = 0$ seraient
équivalents\footnote{« Pour qu'il existe » et « bijection » sont rayés, au
profit de « Alors il existe » et « isomorphisme can. ». L'énoncé dépend du
lemme~1.}.

\pagerange{21}{22}
\subsection*{Matrices et automorphismes (pages 21 et 22)}

\emph{Corollaire 2} (annoncé). Sous les conditions du lemme~1, soient $n
\geqslant 1$, $\mathbb{M}_n$ le faisceau des matrices carrées d'ordre $n$ et
$\mathrm{Gl}_n = \mathbb{M}_n^{*}$, et supposons de plus
\[
  \mathbb{M}_n(\widetilde{B}_0) = \mathbb{M}_n(\widetilde{B}_0)^{*} +
  \mathrm{Im}\bigl(\mathbb{M}_n(B) \to \mathbb{M}_n(\widetilde{B}_0)\bigr)
\]
(ce qui a lieu sous les conditions du N.B. du lemme~1, voir le
lemme~2)\footnote{Cette hypothèse est la note de bas de page de sa main,
appelée par une barre oblique après « lemme~1 ».}. Il y aurait une bijection
respectant les éléments neutres
\[
  \mathrm{Ker}\bigl(H^1(A'/A, \mathrm{Gl}_n) \to H^1(A'_0/A_0,
  \mathrm{Gl}_n)\bigr) \to \mathrm{Ker}\bigl(H^1(A'/A, \mathbb{M}_n) \to
  H^1(A'_0/A_0, \mathbb{M}_n)\bigr) .
\]
Par suite, si $H^1(A'_0/A_0, \mathrm{Gl}_n)$ est trivial et $H^1(A'_0/A_0)
= 0$, on aurait $H^1(A'/A, \mathrm{Gl}_n) \simeq H^1(A'/A, \mathbb{M}_n) =
\mathbb{M}_n\bigl(H^1(A'/A)\bigr)$, et $H^1(A'/A, \mathrm{Gl}_n)$ serait
trivial si et seulement si $H^1(A'/A) = 0$.

La démonstration identifie les deux noyaux à $H^1(I\,C^{\bullet}(A'/A,
\mathbb{M}_n))$ modulo l'action de $\mathrm{Gl}_n(\widetilde{B}_0)$, resp. de
$\mathbb{M}_n(\widetilde{B}_0)$\footnote{Le symbole devant $C^{\bullet}$ est
lu « $Z$ », douteux, aux trois occurrences ; à la page 19 le même complexe
s'écrit nettement $I\,C^{\bullet}(A'/A)$, et c'est ce qu'on lit ici.}. Cette
identification est correcte : c'est la description classique de la fibre
d'un $H^1$ non abélien au-dessus de l'élément neutre, pour une extension à
noyau commutatif. Mais l'action de $\mathrm{Gl}_n(\widetilde{B}_0)$ est
affine — une conjugaison suivie de la translation par un cobord, qui est
une dérivée logarithmique —, et non « induite par celle de
$\mathbb{M}_n(\widetilde{B}_0)$ » comme l'écrit la page : c'est la lacune du
lemme~1 en dimension $n$, aggravée par la conjugaison. La phrase qui devait
conclure s'interrompt sur la note de bas de page\footnote{« Donc il suffit
… en utilisant l'hypothèse », dont trois mots sont illisibles.}.

\emph{Corollaire 3} (annoncé). Même énoncé pour un module libre $E =
A^{(N)}$ ($N$ un ensemble d'indices quelconque), les faisceaux
$\underline{\mathrm{Aut}}(E)$ et $\underline{\mathrm{End}}(E)$, sous
l'hypothèse $\mathrm{End}_{\widetilde{B}_0}(E_{\widetilde{B}_0}) =
\mathrm{Aut}_{\widetilde{B}_0}(E_{\widetilde{B}_0}) +
\mathrm{Im}\,\mathrm{End}_{B}(E_{B})$ ; le noyau du côté additif est
$\bigl[\mathrm{Ker}(H^1(A'/A) \to H^1(A'_0/A_0))^{(N)}\bigr]^N$, puisque
$\mathrm{End}(A^{(N)})$ est le module des matrices à colonnes finies et que
la cohomologie commute aux sommes directes et aux produits\footnote{Les
exposants $(N)$ et $N$ sont repassés à l'encre foncée. Le crochet qui suit,
en partie illisible et dont les crochets ne se referment pas tous, dit que
les données de descente sur $E_{A'}$ triviales sur $E_{A'_0}$ sont alors
effectives et donnent un module libre sur $B$.}. La démonstration est « le
même principe », par les suites exactes de groupes cosimpliciaux $e \to
C^{\bullet}(\underline{\mathrm{Hom}}(E, IE)) \to
C^{\bullet}(\underline{\mathrm{Aut}}(E)) \to
C^{\bullet}_0(\underline{\mathrm{Aut}}(E)) \to e$ et son analogue additif,
qui sont exactes (un endomorphisme qui relève un automorphisme modulo un
idéal de carré nul est un automorphisme) ; elle hérite de la même lacune.

\pagerange{23}{24}
\subsection*{Le lemme 2 : rendre une matrice inversible (pages 23 et 24)}

\emph{Lemme 2.} Soient $B$ un anneau semi-local et $C$ une $B$-algèbre finie,
semi-locale, dont les idéaux maximaux $\mathfrak{n}$ induisent des idéaux
maximaux $\mathfrak{m}$ de $B$. On suppose : (i) les corps
$k(\mathfrak{m})$ sont infinis ; (ii) deux idéaux maximaux distincts de $C$
induisent des idéaux maximaux distincts de $B$, et chaque extension
$k(\mathfrak{n})/k(\mathfrak{m})$ est radicielle (par exemple $C$ fini et
radiciel sur $B$). Alors pour tout $n \geqslant 1$ et tout $u \in
\mathbb{M}_n(C)$, il existe $v \in \mathbb{M}_n(B)$ tel que $u + \varphi(v)
\in \mathrm{Gl}_n(C)$, où $\varphi : B \to C$.

L'inversibilité d'une matrice de $\mathbb{M}_n(C)$ ne dépend que de son image
modulo le radical, et par le lemme chinois on peut prescrire les images de $v$
dans les $\mathbb{M}_n(k(\mathfrak{m}_i))$ ; par (ii), on se ramène à $B = k$
un corps et $C = k(\mathfrak{n})$\footnote{La démonstration de la page 23 est
en grande partie illisible ; le théorème invoqué n'est pas lu. La réduction
donnée ici suit ce qui se lit (« pour $u' \in \mathbb{M}_n(C)$, le fait que
ce soit inversible ne dépend que de l'inversibilité de ses images dans
$\mathbb{M}_n(C/\mathrm{rad}\,C)$ ! »).}. Le polynôme $\det(u + X)$ en les
coefficients de $X$ n'est pas nul (il vaut $1$ en $X = 1 - u$) ; si $k$ est
infini, il ne s'annule pas sur $\mathbb{M}_n(k)$. On observe que
l'hypothèse (i) est en fait superflue sous (ii) : si $k(\mathfrak{m})$ est
fini, il est parfait, l'extension radicielle $k(\mathfrak{n})/k(\mathfrak{m})$
est triviale, et $v = 1 - u$ convient\footnote{L'observation est la nôtre.
La page termine sur un N.B. : « Il faudrait dégager la possibilité d'une
extension résiduelle transcendante $k(t)$, qui fera … en (i) » — sans doute
le procédé classique qui rend les corps résiduels infinis en passant de $B$ à
$B(t)$ ; c'est ce que l'observation rend inutile ici.}.

La page 24 énonce une variante pour un module $E$ libre sur $B$ : pour tout
$v \in \mathrm{End}_C(E_C)$, il existe $u \in \mathrm{End}_B(E)$ tel que $v +
u_C$ soit inversible ; la démonstration, presque entièrement illisible, finit
sur « $u = \mathrm{id}_E - v$ ».

\pagerange{25}{25}
\subsection*{Le théorème (page 25)}

\emph{Théorème.} Soient $A$ un anneau noethérien semi-local de radical
$\mathfrak{m}$ et $A'$ une $A$-algèbre finie\footnote{« ½ local
noethérien » est écrit au-dessus d'« artinien », rayé.}. Posons $A_n =
A/\mathfrak{m}^{n+1}$, $A'_n = A' \otimes_A A_n$ et $B^{(n)} =
H^0(A'_n/A_n)$.

a) Soit $r \geqslant 1$. Les conditions suivantes sont équivalentes :
\begin{itemize}
  \item[(i)] pour tout $n$, $H^1(A'_n/A_n) = 0$ ;
  \item[(ii$_1$)] pour tout $n$, $H^1(A'_n/A_n, \mathbb{G}_m)$ est trivial ;
  \item[(ii$_r$)] pour tout $n$, $H^1(A'_n/A_n, \mathrm{Gl}_r)$ est trivial ;
  \item[(iii)] pour tout $n$, $\mathrm{Spec}\, A'_n \to \mathrm{Spec}\,
  B^{(n)}$ est un morphisme de descente effective pour les modules
  quasi-cohérents plats\footnote{« qu.\ cohér. » et « plats » sont écrits
  l'un au-dessus de l'autre, séparés d'un trait, comme une fraction ; on lit
  « quasi-cohérents plats », ce que la partie c) écrit en toutes lettres. Une
  note diagonale de la marge, presque illisible, semble dire que a) vaut
  pour un anneau local noethérien quelconque.}.
\end{itemize}

b) Ces conditions entraînent, avec $B = H^0(A'/A)$ : (i) $H^1(A'/A) = 0$ ;
(ii) $H^1(A'/A, \mathrm{Gl}_r)$ est trivial pour tout $r$ ; (iii)
$\mathrm{Spec}\, A' \to \mathrm{Spec}\, B$ est un morphisme de descente
effective pour les modules plats \emph{de type fini}\footnote{« de type
fini » est souligné et entouré au crayon.}.

c) Si $A$ est artinien et que les conditions de a) sont remplies, pour tout
$B$-préschéma $X$ plat sur $B$, le morphisme $X' = X \otimes_B A' \to X$ est
un morphisme de descente effective pour les modules quasi-cohérents
plats\footnote{« artinien » est entouré au crayon.}.

Le théorème dit, en somme, qu'un problème de descente le long de $A'/A$ se
décide sur les voisinages infinitésimaux $A_n$, et que sur ceux-ci la
descente des modules plats, celle des fibrés de rang $r$ et l'annulation
d'un $H^1$ additif sont une seule et même condition. Pour $A$ artinien, $A_n
= A$ dès que $n$ est grand, et b), c) ne demandent pas de passage à la
limite ; pour $A$ complet, la partie b) est le passage aux « modules
formels » du titre.

\pagerange{26}{29}
\subsection*{La démonstration (pages 26 à 29)}

\emph{a).} L'équivalence de (i) et (ii$_r$) est « immédiate par récurrence
sur $n$ » par les lemmes~1 et~2 et le corollaire~2 : pour $n = 0$, $A_0 =
A/\mathfrak{m}$ est un produit de corps, $A'_0$ est plat sur $A_0$, et la
descente fidèlement plate (sur le support de $A'_0$) donne $H^1 = 0$ et
$H^1(\mathrm{Gl}_r)$ trivial ; on passe de $n-1$ à $n$ par l'idéal
$\mathfrak{m}^n/\mathfrak{m}^{n+1}$, de carré nul, l'hypothèse du lemme~1
étant remplie par son N.B. et celle du corollaire~2 par le lemme~2. La
récurrence est correcte \emph{si} les énoncés du lemme~1 et du corollaire~2
le sont, ce que, on l'a vu, le dossier ne démontre pas\footnote{La page
écrit l'hypothèse vérifiée sous la forme $\mathbb{M}_r(B^{(n)}) =
\mathbb{M}_r(B^{(n)})^{*} + \mathrm{Im}\,\mathbb{M}_r(B^{(n)})$, très serrée,
dont le dernier terme est douteux ; elle invoque « le lemme~2 (ii) » et « la
théorie des anneaux artiniens ».}.

Ensuite, (ii$_r$) dit que tout $A'_n$-module libre de rang $r$ muni d'une
donnée de descente relativement à $A'_n/B^{(n)}$ est effectif, de module
descendu libre : c'est la descente effective pour les modules libres de rang
fini\footnote{« rang fini » en interligne au-dessus de « type fini »,
rayé.}. Donc (iii) entraîne (ii$_r$). Pour (i) $\Rightarrow$ (iii) : (i) est
stable par changement de base plat et (iii) est invariant par changement de
base fidèlement plat, ce qui ramène au cas des extensions résiduelles
triviales, puis à $A = A_n = B^{(n)}$ local ; un $A'_n$-module plat est alors
libre sur chaque composante locale, et l'effectivité résulte du
corollaire~3 du lemme~1 et du lemme~2\footnote{Une grande partie des pages 26
et 27 est illisible, dont la phrase « Je dis que les modules descendus … » et
la justification de la liberté. La page identifie entre crochets la donnée de
descente relativement à $A'_n/A_n$ à celle relativement à $A'_n/B^{(n)}$ ;
c'est exact, et la justification est la nôtre : pour $b \in B^{(n)}$ on a déjà
$b \otimes 1 = 1 \otimes b$ dans $A'_n \otimes_{A_n} A'_n$, de sorte que les
produits tensoriels sur $A_n$ et sur $B^{(n)}$ coïncident, et les complexes
d'Amitsur aussi.}.

\emph{b) (i).} La formation de $H^1(A'/A)$ commute au changement de base
plat, et $A \to \widehat{A}$ est fidèlement plat ; on peut supposer $A$
complet. Alors $C^{\bullet}(A'/A) = \varprojlim_n C^{\bullet}(A'_n/A_n)$,
les termes étant des $A$-modules finis. Les systèmes de cochaînes ont des
transitions surjectives, et le système des $B^{(n)} = H^0(A'_n/A_n)$,
modules de longueur finie, vérifie la condition de Mittag-Leffler ; donc
$\varprojlim^1 B^{(n)} = 0$ et
\[
  H^1(A'/A) = \varprojlim_n H^1(A'_n/A_n) = 0 ,
\]
ce que la page résume par « en vertu de M.L. »\footnote{La page écrit seulement « en vertu de M.L. » ; la condition sur
les $H^0$, qui est celle dont la suite exacte de Milnor a besoin, est
précisée par nous. La page ajoute qu'on peut aussi déduire formellement (i)
de (iii) « en faisant $m = 1$, et $m = 2$ », sans le faire.}.

\emph{b) (ii)} résulte de (iii) « par le même argument que dans a) ».

\emph{b) (iii).} On peut supposer $A$ complet, $A = B$ et $A$ local. Soit
$E'$ un $A'$-module plat de type fini muni d'une donnée de descente, $E =
H^0_D(A'/A, E')$ ; il faut montrer que $E$ est plat et que $E \otimes_A A'
\to E'$ est un isomorphisme. Comme $\varprojlim$ est exact à gauche, $E =
\varprojlim_n E^{(n)}$ avec $E^{(n)} = H^0_D(A'_n/A_n, E'_n)$, et par a)
(iii), $E^{(n)}$ est un $B^{(n)}$-module plat de type fini avec $E^{(n)}
\otimes_{B^{(n)}} A'_n \simeq E'_n$. La page conclut par « un lemme facile de
passage à la limite » que $E$ est plat sur $A = \varprojlim B^{(n)}$, et que
$E \otimes_A A' \to E'$, limite des isomorphismes $E^{(n)} \otimes A'_n \to
E'_n$, est un isomorphisme\footnote{Le lemme n'est pas écrit. Il n'est pas
tout à fait immédiat : le système $(B^{(n)})$ n'est pas le système adique
$(B/\mathfrak{m}^{n+1}B)$, et il faut savoir qu'il lui est
pro-isomorphe, ce qui découle du lemme d'Artin–Rees appliqué au noyau de
$A' \rightrightarrows A' \otimes_A A'$ ; on le signale sans le
démontrer.}.

\emph{c)} est seulement entamé : il s'agit de montrer que $X' \to X$ est de
descente, puis de descente effective, en se ramenant à $A = B$ et aux $A_n$,
avec l'hypothèse a) sous la forme (iii) ; la page s'arrête dans un
crochet\footnote{Au bas de la page 29, isolée et sans texte, la ligne $A \to
A' \rightrightarrows A''$.}.

\pagerange{30}{31}
\section*{IV. Les deux derniers feuillets (pages 30 et 31)}

La page 30 est une feuille brune vierge, sans doute la chemise de dos de la
liasse. La page 31 est d'une autre main, propre et en anglais, sans
intervention repérée de Grothendieck. Elle porte sur les idéaux
d'équivalence de $A' \otimes_k A'$ pour l'algèbre $A' = k \oplus V$ de carré
nul, sous la condition de transitivité $p_{13}(I) \subset p_{12}(I) +
p_{23}(I)$ : c'est le cas que la page 18 laisse en suspens. Le titre de
l'inventaire, « Cf Notes Murre-Levelt », invite à y voir une feuille de ces
notes ; rien sur la feuille ne nomme son auteur, et l'on n'en dit pas
davantage\footnote{La transcription en résume le contenu dans une note ; on ne
le reprend pas ici.}.

\end{document}
